% Sezione 4: Modelli Cinematici e Controllo
\section{Modellazione Cinematica e Controllo}

\begin{frame}
    \frametitle{Modello a Biciclo Sterzante}
    \textbf{Semplificazione}: ruote sullo stesso asse $\rightarrow$ ruota virtuale unica per asse

    \vspace{0.3cm}

    \textbf{Equazioni Cinematiche}
    \begin{equation*}
        \begin{cases}
            \dot{x}      = v \cos(\theta+\beta) \\
            \dot{y}      = v \sin(\theta+\beta) \\
            \dot{\theta} = \frac{v\cos(\beta)}{l_f+l_r}\left(\tan(\delta_f)-\tan(\delta_r)\right)
        \end{cases}
    \end{equation*}

    dove:
    \begin{align*}
        v     & =  \frac{r(\omega_f \cos\delta_f + \omega_r \cos\delta_r)}{2\cos(\beta)}                                    \\
        \beta & =  \tan^{-1}\left(\frac{l_f \tan(\delta_r) + l_r \tan(\delta_f)}{l_f + l_r}\right) \quad \text{(side-slip)}
    \end{align*}

    \vspace{0.2cm}
\end{frame}
\begin{frame}
    \begin{block}{Vantaggi}
        \begin{itemize}
            \item Testato in letteratura.
            \item Modello dinamico disponibile e controllato con successo con LQR.
        \end{itemize}
    \end{block}
    \begin{alertblock}{Svantaggi: Singolarità}
        \begin{itemize}
            \item $\beta \rightarrow \infty$ se $\delta_f$ o $\delta_r \rightarrow \pi/2$ (tangente infinita)
            \item \textbf{Crab/Spin impossibile}: se $v=0$ allora $\omega_f = \omega_r = 0 \Rightarrow \dot{\theta}=0$
        \end{itemize}
    \end{alertblock}
\end{frame}

\begin{frame}
    \frametitle{Modello a Doppia Ruota Sterzante: Vincoli Non-Olonomi}
    \textbf{Vincolo di non slittamento} per ruota $i \in \{f, r\}$:

    \vspace{0.2cm}

    Velocità punto contatto ruota nel body frame:
    \begin{equation*}
        \mathbf{v}_i^b = \begin{bmatrix}
            \dot{x}_c \mp \dot{\theta}\frac{W}{2} \\
            \dot{y}_c \pm \dot{\theta}\frac{L}{2}
        \end{bmatrix}
    \end{equation*}

    \textbf{Versori}: normale $\mathbf{n}_i = \begin{bmatrix}
            -\sin(\delta_i) \\
            \cos(\delta_i)
        \end{bmatrix}$, parallelo $\mathbf{u}_i = \begin{bmatrix}
            \cos(\delta_i) \\
            \sin(\delta_i)
        \end{bmatrix}$

    \textbf{Condizione di non slittamento}: $\mathbf{n}_i^T \mathbf{v}_i^b = 0$
    \begin{equation*}
        -\sin(\delta_i)\left(\dot{x}_c \mp \dot{\theta}\frac{W}{2}\right) + \cos(\delta_i)\left(\dot{y}_c \pm \dot{\theta}\frac{L}{2}\right) = 0
    \end{equation*}

    \vspace{0.1cm}

    \textbf{Vincolo di puro rotolamento}: $\mathbf{u}_i^T \mathbf{v}_i^b = r\omega_i$
    \begin{equation*}
        \cos(\delta_i)\left(\dot{x}_c \mp \dot{\theta}\frac{W}{2}\right) + \sin(\delta_i)\left(\dot{y}_c \pm \dot{\theta}\frac{L}{2}\right) = r\omega_i
    \end{equation*}

    \vspace{0.2cm}

    \small{$(f): L/2, W/2$; $(r): -L/2, -W/2$}
\end{frame}

\begin{frame}
    \frametitle{Cinematica Diretta: Matrice dei Vincoli Estesa}
    \textbf{Sistema sovradeterminato} (4 equazioni, 3 incognite):

    \begin{equation*}
        \mathbf{A}_{ext}^T \dot{\mathbf{q}} = \mathbf{b}
    \end{equation*}

    dove:
    \begin{equation*}
        \mathbf{A}_{ext}^T = \begin{bmatrix}
            -\sin(\delta_f) & \cos(\delta_f) & \frac{W}{2}\sin(\delta_f) - \frac{L}{2}\cos(\delta_f)  \\
            -\sin(\delta_r) & \cos(\delta_r) & -\frac{W}{2}\sin(\delta_r) + \frac{L}{2}\cos(\delta_r) \\
            \cos(\delta_f)  & \sin(\delta_f) & -\frac{W}{2}\cos(\delta_f) - \frac{L}{2}\sin(\delta_f) \\
            \cos(\delta_r)  & \sin(\delta_r) & \frac{W}{2}\cos(\delta_r) - \frac{L}{2}\sin(\delta_r)
        \end{bmatrix}
        , \quad
        \mathbf{b} = \begin{bmatrix}
            0 \\ 0 \\ r\omega_f \\ r\omega_r
        \end{bmatrix}
    \end{equation*}

    \textbf{Soluzione ai minimi quadrati} (pseudoinversa di Moore-Penrose):
    \begin{equation*}
        \dot{\mathbf{q}} = \mathbf{A}_{ext}^\dagger \mathbf{b} = \left(\mathbf{A}_{ext}^T \mathbf{A}_{ext}\right)^{-1} \mathbf{A}_{ext}^T \mathbf{b}
    \end{equation*}
\end{frame}

\begin{frame}
    \frametitle{Cinematica Diretta: Forma Compatta}
    \textbf{Soluzione finale} nel body frame:

    \begin{equation*}
        \begin{bmatrix}
            \dot{x}_c \\ \dot{y}_c \\ \dot{\theta}
        \end{bmatrix}
        = \frac{r}{2}
        \begin{bmatrix}
            \cos(\delta_f) & \cos(\delta_r) \\
            \sin(\delta_f) & \sin(\delta_r) \\
            a_{13}         & a_{23}
        \end{bmatrix}
        \begin{bmatrix}
            \omega_f \\ \omega_r
        \end{bmatrix}
    \end{equation*}

    dove i coefficienti $a_{13}$, $a_{23}$ dipendono da $\delta_f$, $\delta_r$, $L$, $W$:
    \begin{align*}
        a_{13} & = \frac{2(L\sin(\delta_f) - W\cos(\delta_f))}{L^2 + W^2} \\
        a_{23} & = \frac{2(W\cos(\delta_r) - L\sin(\delta_r))}{L^2 + W^2}
    \end{align*}
\end{frame}
\begin{frame}
    \begin{block}{Vantaggi}
        \begin{itemize}
            \item Trasformabile facilmente in frame globale e viceversa (pre-mnoltiplicando $R(\theta)$/$R^T(\theta)$)
            \item \textbf{Nessuna singolarità cinematica} per qualsiasi $\delta_f, \delta_r$ (Pseudoinversa mai singolare)
            \item Gestisce \textbf{crab} ($\dot{x}_c=0$, $\dot{y}_c\neq0$) e \textbf{spin} ($\dot{x}_c=\dot{y}_c=0$, $\dot{\theta}\neq0$)
        \end{itemize}
    \end{block}
    \begin{alertblock}{Svantaggi}
        \begin{itemize}
            \item No modellazione dinamica (forze, inerzie)
            \item No modellazione ruote libere
            \item Soluzione ai minimi quadrati: possibilità di rompere vincoli non-olonomi
        \end{itemize}
    \end{alertblock}
\end{frame}

\begin{frame}
    \frametitle{Centro Istantaneo di Rotazione (ICR)}
    \textbf{Concetto}: punto attorno al quale il veicolo sta ruotando in un dato istante

    \vspace{0.3cm}

    L'ICR si trova all'\textbf{intersezione degli assi di rotazione delle ruote} (perpendicolari alla direzione di avanzamento):

    \begin{equation}
        \begin{bmatrix}
            x_{ICR} \\ y_{ICR}
        \end{bmatrix}
        = \mathbf{p}_i^b + s_i \begin{bmatrix}
            -\sin(\delta_i) \\ \cos(\delta_i)
        \end{bmatrix}
        \quad i \in \{f, r\}
    \end{equation}

    \textbf{Distanza da ICR}: $R_i = |s_i|$ (raggio di curvatura per ruota $i$)

    \vspace{0.2cm}

    \textbf{Condizione di compatibilità} (moto senza slittamento):
    \begin{equation}
        \dot{\theta} R_i = r\omega_i \quad \Rightarrow \quad R_r \omega_f = R_f \omega_r
    \end{equation}

    \small{Caso $\delta_f \parallel \delta_r$: ICR $\to \infty$, traslazione pura con $\omega_f = \omega_r$}
\end{frame}

\begin{frame}
    \frametitle{Cinematica Inversa: Modello a Biciclo}
    \textbf{Problema}: dato $(\dot{x}_d, \dot{y}_d, \dot{\theta}_d)$, calcolare $(\delta_f, \delta_r, \omega_f, \omega_r)$

    \vspace{0.3cm}

    \textbf{Passo 1}: Calcolare velocità e side-slip desiderati
    \begin{align}
        v_d     & = \sqrt{\dot{x}_d^2 + \dot{y}_d^2} \nonumber       \\
        \beta_d & = \atanyx{\dot{y}_d}{\dot{x}_d} - \theta \nonumber
    \end{align}

    \textbf{Passo 2}: Risolvere sistema per angoli di sterzo (definendo $S = \tan(\beta_d)$, $D = \frac{l\dot{\theta}_d}{v_d \cos(\beta_d)}$)
    \begin{align}
        \delta_f & = \arctan\left(S + \frac{l_f}{l} D\right) \\
        \delta_r & = \arctan\left(S - \frac{l_r}{l} D\right)
    \end{align}

    \textbf{Passo 3}: Calcolare $\omega_f$, $\omega_r$ da raggi ICR $R_f$, $R_r$ e condizione compatibilità
\end{frame}

\begin{frame}
    \frametitle{Cinematica Inversa Biciclo: Caso $\dot{\theta}_d = 0$ (Parte 1)}
    \textbf{Derivazione per traslazione pura}

    \vspace{0.3cm}

    Quando $\dot{\theta}_d = 0$, l'ICR $\to \infty$ e le direzioni di avanzamento delle ruote diventano parallele.

    \vspace{0.3cm}

    \textbf{Step 1}: Imporre $\delta_f = \delta_r = \delta$ (sterzi paralleli)

    \vspace{0.3cm}

    \textbf{Step 2}: Dalla relazione del side-slip angle:
    \begin{equation}
        \beta_d = \tan^{-1}\left(\frac{l_f + l_r}{l_f + l_r}\tan(\delta)\right) = \delta
    \end{equation}

    \vspace{0.3cm}

    Quindi il side-slip angle coincide con l'angolo di sterzo quando le ruote sono parallele.
\end{frame}

\begin{frame}
    \frametitle{Cinematica Inversa Biciclo: Caso $\dot{\theta}_d = 0$ (Parte 2)}
    \textbf{Calcolo angoli di sterzo e velocità}

    \vspace{0.3cm}

    \textbf{Step 3}: Combinando $\beta_d = \delta$ con $\beta_d = \atanyx{\dot{y}_d}{\dot{x}_d} - \theta$:
    \begin{equation}
        \delta_f = \delta_r = \delta = \atanyx{\dot{y}_d}{\dot{x}_d} - \theta
    \end{equation}

    \vspace{0.3cm}

    \textbf{Step 4}: Le velocità angolari sono uguali (condizione di compatibilità per ICR $\to \infty$):
    \begin{equation}
        \omega_f = \omega_r = \frac{v_d}{r} = \frac{\sqrt{\dot{x}_d^2 + \dot{y}_d^2}}{r}
    \end{equation}
\end{frame}

\begin{frame}
    \frametitle{Cinematica Inversa: Modello Bisteer}
    \textbf{Soluzione diretta dai vincoli di non-slittamento}:

    \vspace{0.3cm}

    \textbf{Angoli di sterzo} (dalle eq. non-slip):
    \begin{align}
        \delta_f & = \atanyx{\dot{y}_d + \dot{\theta}_d \frac{L}{2}}{\dot{x}_d - \dot{\theta}_d \frac{W}{2}} \\
        \delta_r & = \atanyx{\dot{y}_d - \dot{\theta}_d \frac{L}{2}}{\dot{x}_d + \dot{\theta}_d \frac{W}{2}}
    \end{align}

    \textbf{Velocità angolari ruote} (dalle eq. rotolamento):
    \begin{align}
        \omega_f & = \frac{(\dot{y}_d + \dot{\theta}_d \frac{L}{2}) \sin(\delta_f) + (\dot{x}_d - \dot{\theta}_d \frac{W}{2}) \cos(\delta_f)}{r} \\
        \omega_r & = \frac{(\dot{y}_d - \dot{\theta}_d \frac{L}{2}) \sin(\delta_r) + (\dot{x}_d + \dot{\theta}_d \frac{W}{2}) \cos(\delta_r)}{r}
    \end{align}

    \textbf{Caso speciale $\dot{\theta}_d = 0$}: $\delta_f = \delta_r = \atanyx{\dot{y}_d}{\dot{x}_d}$, $\omega_f = \omega_r = \frac{v_d}{r}$
\end{frame}

\begin{frame}
    \frametitle{Setup Simulazioni}
    \textbf{Parametri Robot e Attuatori}
    \begin{itemize}
        \item Dimensioni: $L = 2$ m, $W = 1$ m, $r = 0.198$ m
        \item Dinamica attuatori: modello 1$^\circ$ ordine $P(s) = \frac{K}{\tau s + 1}$
              \begin{itemize}
                  \item Trazione: $K_{\omega} = 1$, $\tau_\omega = 1$ s, $\omega_{max} = 10$ rad/s
                  \item Sterzo: $K_{\delta} = 1$, $\tau_\delta = 2$ s, $\delta_{max} = \pm 120^\circ$
              \end{itemize}
    \end{itemize}

    \vspace{0.3cm}

    \textbf{Traiettorie di Test}
    \begin{enumerate}
        \item \textbf{Straight}: linea retta verso Est ($v_{ref} = 0.5$ m/s, $T = 5$ s)
        \item \textbf{Arc}: arco circolare ($R = 3$ m, $v_{ref} = 0.6$ m/s, $\omega = 0.2$ rad/s, $T = 8$ s)
        \item \textbf{Crab}: traslazione laterale Nord ($v_{ref} = 0.4$ m/s, $T = 5$ s)
        \item \textbf{Spin}: rotazione sul posto ($\omega_{ref} = \pi/6$ rad/s, $T = 6$ s)
    \end{enumerate}

    \vspace{0.2cm}
    \small{Metrica: RMSE posizione e orientamento vs. traiettoria desiderata}
\end{frame}

\begin{frame}
    \frametitle{Schema di Controllo: Feedforward (Catena Aperta)}
    \begin{figure}
        \centering
        \includegraphics[width=0.95\textwidth]{4/FF_scheme.pdf}
    \end{figure}

    \textbf{Caratteristiche}:
    \begin{itemize}
        \item Nessun feedback dallo stato misurato
        \item Cinematica inversa calcola input da traiettoria desiderata
        \item Dinamica attuatori introduce ritardi non compensati
        \item Errori si accumulano senza correzione
    \end{itemize}
\end{frame}

\begin{frame}
    \frametitle{Risultati Feedforward: Traiettoria Arc}
    \begin{figure}
        \centering
        \includegraphics[width=0.85\textwidth]{4/curved_east_to_north_feedforward_bicycle.png}
        \caption{Modello Bicycle - Traiettoria circolare (Feedforward)}
    \end{figure}
\end{frame}

\begin{frame}
    \frametitle{Risultati Feedforward: Traiettoria Arc}
    \begin{figure}
        \centering
        \includegraphics[width=0.85\textwidth]{4/curved_east_to_north_feedforward_bisteer.png}
        \caption{Modello Bisteer - Traiettoria circolare (Feedforward)}
    \end{figure}
\end{frame}

\begin{frame}
    \frametitle{Risultati Feedforward: Confronto e Analisi}
    \begin{columns}
        \column{0.5\textwidth}
        \begin{table}
            \centering
            \tiny
            \begin{tabular}{@{}lcccc@{}}
                \toprule
                \textbf{Traiett.} & \textbf{Mod.} & \textbf{Supp.} & \textbf{Pos. [mm]}                & \textbf{Orient. [$^\circ$]}      \\
                \midrule
                \multirow{2}{*}{Str.}
                                  & Bic.          & x              & 418.05                            & 0.00                             \\
                                  & Bis.          & x              & 418.05                            & 0.00                             \\
                \midrule
                \multirow{2}{*}{Arc}
                                  & Bic.          & x              & 749.92                            & 20.80                            \\
                                  & Bis.          & x              & \textbf{694.70}                   & \textbf{18.85}                   \\
                \midrule
                \multirow{2}{*}{Cra.}
                                  & Bic.          & --             & \textcolor{red}{\textbf{1154.82}} & 0.00                             \\
                                  & Bis.          & x              & \textbf{334.44}                   & 0.00                             \\
                \midrule
                \multirow{2}{*}{Spi.}
                                  & Bic.          & --             & 0.00                              & \textcolor{red}{\textbf{103.93}} \\
                                  & Bis.          & x              & 0.00                              & \textbf{25.93}                   \\
                \bottomrule
            \end{tabular}
        \end{table}

        \column{0.45\textwidth}
        \textbf{Osservazioni}:
        \begin{itemize}
            \item Bicycle FALLISCE su Crab e Spin (singolarità)
            \item Bisteer gestisce tutte le manovre
            \item Errori elevati per assenza feedback
        \end{itemize}
    \end{columns}
\end{frame}

\begin{frame}
    \frametitle{Schema di Controllo: PID + Feedforward}
    \begin{figure}
        \centering
        \includegraphics[width=\textwidth]{4/FF_PID.pdf}
    \end{figure}
\end{frame}
\begin{frame}
    \frametitle{Dettagli Implementazione PID}
    \textbf{Legge di controllo}:
    \begin{align}
        e^b           & = R^T(\theta) \begin{bmatrix} x_d - x \\ y_d - y \\ \theta_d - \theta \end{bmatrix} \quad \text{(errore in body frame)} \nonumber \\
        u_i(t)        & = K_{p,i} e_i^b + K_{i,i} \int e_i^b \, dt + K_{d,i} \frac{d e_i^b}{dt}, \quad i \in \{x,y,\theta\} \nonumber                     \\
        \dot{q}_{cmd} & = \dot{q}_d + u \quad \text{(velocità corrette in global frame)} \nonumber
    \end{align}
\end{frame}

\begin{frame}
    \frametitle{Dettagli Implementazione PID}
    \textbf{Guadagni utilizzati}:
    \begin{table}
        \centering
        \small
        \begin{tabular}{cccccc}
            \toprule
            \textbf{Comp.} & $K_p$ & $K_i$ & $K_d$ & $I_{max}$ & $N$ \\
            \midrule
            $x, y$         & 0.3   & 0.05  & 0.1   & 2         & 20  \\
            $\theta$       & 0.2   & 0.02  & 0.05  & 1         & 20  \\
            \bottomrule
        \end{tabular}
    \end{table}

    \vspace{0.3cm}

    \textbf{Accorgimenti implementativi}:
    \begin{itemize}
        \item \textbf{Derivata filtrata}: $D(s) = \frac{K_d s}{s/N + 1}$ (riduce rumore)
        \item \textbf{Anti-windup}: saturazione integrale per evitare wind-up
        \item \textbf{Scaling proporzionale}: mantiene rapporto tra $u_x, u_y$ sotto $v_{max}$
        \item \textbf{Saturazione angolare}: $|u_\theta| \leq \omega_{max}$
    \end{itemize}
\end{frame}

\begin{frame}
    \frametitle{Risultati PID+FF: Traiettoria Arc}
    \begin{figure}
        \centering
        \includegraphics[width=0.85\textwidth]{4/curved_east_to_north_pid_bicycle.png}
        \caption{Modello Bicycle - Traiettoria circolare (PID+FF)}
    \end{figure}
\end{frame}

\begin{frame}
    \frametitle{Risultati PID+FF: Traiettoria Arc}
    \begin{figure}
        \centering
        \includegraphics[width=0.85\textwidth]{4/curved_east_to_north_pid_bisteer.png}
        \caption{Modello Bisteer - Traiettoria circolare (PID+FF)}
    \end{figure}
\end{frame}

\begin{frame}
    \frametitle{Risultati PID+FF: Confronto e Miglioramenti}
    \begin{columns}
        \column{0.55\textwidth}
        \begin{table}
            \centering
            \tiny
            \begin{tabular}{@{}lcccc@{}}
                \toprule
                \textbf{Traiett.} & \textbf{Mod.} & \textbf{Supp.} & \textbf{Pos. [mm]}                & \textbf{Orient. [$^\circ$]}      \\
                \midrule
                \multirow{2}{*}{Str.}
                                  & Bic.          & x              & 233.96                            & 0.00                             \\
                                  & Bis.          & x              & 233.96                            & 0.00                             \\
                \midrule
                \multirow{2}{*}{Arc}
                                  & Bic.          & x              & 231.27                            & 12.53                            \\
                                  & Bis.          & x              & \textbf{229.57}                   & \textbf{11.66}                   \\
                \midrule
                \multirow{2}{*}{Cra.}
                                  & Bic.          & --             & \textcolor{red}{\textbf{1154.82}} & 0.00                             \\
                                  & Bis.          & x              & \textbf{187.17}                   & 0.00                             \\
                \midrule
                \multirow{2}{*}{Spi.}
                                  & Bic.          & --             & 0.00                              & \textcolor{red}{\textbf{103.93}} \\
                                  & Bis.          & x              & 0.00                              & \textbf{16.14}                   \\
                \bottomrule
            \end{tabular}
        \end{table}

        \column{0.45\textwidth}
        \textbf{Miglioramenti vs FF}:
        \begin{itemize}
            \item Str: $418 \to 234$ mm ($-44\%$)
            \item Arc: $695 \to 230$ mm ($-67\%$)
            \item Crab: $334 \to 187$ mm ($-44\%$)
            \item Spin: $26^\circ \to 16^\circ$ ($-38\%$)
        \end{itemize}
    \end{columns}
\end{frame}

\begin{frame}
    \frametitle{Conclusioni: Confronto Modelli Cinematici}
    \begin{columns}
        \column{0.5\textwidth}
        \textbf{Modello Biciclo}
        \begin{itemize}
            \item[+] Testato in letteratura
            \item[+] Estensione dinamica disponibile
            \item[-] \textbf{Singolarità}: Crab e Spin impossibili
            \item[-] $\beta \to \infty$ per $\delta \to \pi/2$
        \end{itemize}

        \column{0.5\textwidth}
        \textbf{Modello Bisteer}
        \begin{itemize}
            \item[+] \textbf{Nessuna singolarità}
            \item[+] Supporta tutte le manovre
            \item[+] Prestazioni migliori con PID
            \item[!] Pseudoinversa può violare vincoli
        \end{itemize}
    \end{columns}

    \vspace{0.5cm}

    \textbf{Scelta finale}: \textcolor{green!60!black}{\textbf{Modello Bisteer}} per maggiore versatilità e robustezza
\end{frame}
